\documentclass{article}
\usepackage[UTF8]{ctex}
\usepackage{amsmath,amssymb,amsfonts}
\usepackage{enumitem}
\usepackage{geometry}
\geometry{a4paper, margin=1in}
% === 错题本排版优化（不参与题目标识） ===
\setlength{\parindent}{0pt}
\setlength{\parskip}{0.35em plus 0.10em minus 0.08em}
\setlength{\emergencystretch}{3em}
\setlist{itemsep=0.25em, topsep=0.40em, parsep=0pt}
\allowdisplaybreaks[2]
% === 错题本排版优化结束 ===
\begin{document}
\section*{哈尔滨工业大学2026年数学分析真题}
\section*{83．哈尔滨工业大学2026年数学分析试卷}

\subsection{题目1}
1．函数 \(f(x)\) 在开区间 \((a, b)\) 与 \((b, c)\) 均一致连续，则 \(f(x)\) 在 \((a, b) \cup(b, c)\) 上一致连续．

\subsection{题目2}
2．函数 \(f(x)\) 在 \([a, b]\) 上可导，且在某点 \(c \in(a, b)\) 处 \(f^{\prime}(c)>0\) ，则存在邻域 \((c-\delta, c+\delta)\) 使得 \(f(x)\)在该邻域单调递增。

\subsection{题目3}
3．设函数 \(F(x)\) 在 \([a, b]\) 上可导且 \(F^{\prime}(x)=f(x), x \in[a, b]\) ，则 \(\int_{a}^{b} f(x) \mathrm{d} x=F(b)-F(a)\) ．

\subsection{题目4}
4．设正项级数 \(\sum_{n=1}^{\infty} a_{n}\) 收敛，则 \(\sum_{n=1}^{\infty} \frac{\sqrt{a_{n}}}{\sqrt{n}}\) 收敛。

\subsection{题目5}
5．设函数 \(f(x, y)\) 在 \((0,0)\) 处的所有方向极限都存在且相等，则 \(\lim _{(x, y) \rightarrow(0,0)} f(x, y)\) 存在．
\end{document}
